\subsection{The counter-evaluation equality rule}

Reflexivity and the single counter-evaluation rule are
\[
 \frac{}{t=t},\qquad
 \frac{s=t\qquad R(Ss,a;i,u.h)=r}
      {r=h[t/i,R(t,a;i,u.h)/u]}.
\tag{CE}
\]
Every term in this display denotes its canonical form, and body substitutions
are simultaneous and capture avoiding. In particular the second occurrence
of $R$ becomes $Sa$ when $t$ is literally zero. The seed and update body
are read from the second premise and are not rewritten at passive occurrences
of $s$.

Two applications recover common-left equality composition. Put
$K_s=R(Ss,0;i,u.i)$. Given $p:s=t$ and $q:s=r$, apply CE first to $q$
and $K_s=K_s$, obtaining $K_s=r$, and then to $p$ and that result,
obtaining $r=t$. Therefore symmetry follows from $s=s$ and $s=t$,
and transitivity follows by reversing the first premise and joining at its
common left endpoint. No induction or logical operation is used.

For any fixed numerical template $f(v)$, put
$L_s=R(Ss,0;i,u.f(i))$. CE on $s=s$ and $L_s=L_s$ gives
$L_s=f(s)$. Applying CE again with $s=t$ gives $f(s)=f(t)$.
This derives unrestricted term congruence, including repeated occurrences
and occurrences under nested recursor binders. Joining this equation with
$f(s)=r$ derives crossed transport $r=f(t)$. Finally, CE on the two
appropriate reflexivities is exactly the successor computation equation stated by CE.

\subsection{Predecessor and every ordinary primitive recursion}

Define the following finite terms, with all auxiliary binders fresh:
\[
\begin{aligned}
 P(n)&=R(n,0;i,u.i),\\
 C(n,a,b)&=P(R(n,a;i,u.Sb)),\\
 \operatorname{Rec}(n,a;i,u.h)
   &=P\bigl(R(n,a;i,v.S(h(i,P(v))))\bigr).
\end{aligned}
\]
Computation and congruence alone prove
\[
\begin{aligned}
 P(0)&=1,& P(Sn)&=n,\\
 C(0,a,b)&=a,& C(Sn,a,b)&=b,\\
 \operatorname{Rec}(0,a;h)&=a,&
 \operatorname{Rec}(Sn,a;h)&=h(n,\operatorname{Rec}(n,a;h)).
\end{aligned}
\tag{Compiler equations}
\]
Indeed the inner trajectory
$T(n)=R(n,a;i,v.S(h(i,P(v))))$ has $T(0)=Sa$ and
$T(Sn)=S(h(n,P(T(n))))$. Applying $P(Sx)=x$ gives both equations.
This calculation explains why the successor seed does not restrict allowed
initial states: the state is stored as a successor, and the outer $P$
removes exactly that wrapper after each observation. No informal claim
about a trajectory's values substitutes for the displayed equations.

The term $P$ is not the usual predecessor at zero. Define instead
\[
 p(n)=C(n,0,P(n)),\qquad \operatorname{sg}(n)=C(n,0,1).
\]
Then $p(0)=0$, $p(Sn)=n$, $\operatorname{sg}(0)=0$, and
$\operatorname{sg}(Sn)=1$. Applying $P$ to $Sx=Sy$ proves
successor injectivity. Applying $\operatorname{sg}$ to $Sx=0$ gives
$1=0$, contradicting the ordinary-realization separation axiom
$0=1\to\bot$.

Scalar iteration is the special case
\[
 I(n,a;u.h)=\operatorname{Rec}(n,a;i,u.h),
 \qquad i\notin\operatorname{FV}(h).
\]
It has $I(0,a;h)=a$ and $I(Sn,a;h)=h(I(n,a;h))$ for every seed $a$.
This recovers either raw scalar iteration up to provable equality or its
zero-normal version under normalizing substitution. The counter-aware
language can compile general primitive recursion directly: zero,
successor, and projections have their evident finite terms; composition is
substitution; recursion is the displayed $\operatorname{Rec}$ construction.
There is no need to construct pairing before compiling recursion.
This is an external compiler for each fixed primitive-recursive description,
not an internally total universal evaluator for description codes.

